\documentclass[12pt,a4paper]{article}

\usepackage[utf8]{inputenc}
\usepackage[T1]{fontenc}
\usepackage{geometry}
\usepackage{setspace}
\usepackage{amsmath,amssymb,amsfonts,bm}
\usepackage{graphicx}
\usepackage{booktabs}
\usepackage{array}
\usepackage{enumitem}
\usepackage{caption}
\usepackage{subcaption}
\usepackage{hyperref}
\usepackage{float}
\usepackage{longtable}
\usepackage{multirow}
\usepackage{titlesec}
\usepackage{lmodern}
\usepackage{indentfirst}

\geometry{margin=1in}
\setstretch{1.15}
\setlength{\parindent}{1.25cm}
\setlength{\parskip}{0pt}
\setlist[itemize]{leftmargin=2em}
\hypersetup{hidelinks}

% Resolve figures directly from the analysis outputs of the project.
\graphicspath{%
    {results/analysis/}%
    {results/analysis/flyover/}%
    {results/analysis/av_predictions_2020/}%
    {results/perception_logs/day7_thesis_package/}%
}

% Figure command: #1 image (resolved via \graphicspath), #2 caption, #3 small
% description, #4 label. \includegraphics searches \graphicspath, so the figures
% in results/analysis and its subfolders resolve from the project root.
\newcommand{\figorbox}[4]{%
    \begin{figure}[H]
        \centering
        \includegraphics[width=0.90\textwidth]{#1}
        \caption{#2\\\small #3}\label{#4}
    \end{figure}
    \vspace{0.3em}
}

\begin{document}

\hypersetup{pageanchor=false}
\begin{titlepage}
    \thispagestyle{empty}
    \begin{center}
        \vspace*{1.45cm}

        {\large THE MINISTRY OF SCIENCE AND HIGHER EDUCATION\par}
        \vspace{0.22cm}
        {\large OF THE RUSSIAN FEDERATION\par}

        \vspace{1.25cm}

        {\large ITMO University\par}
        \vspace{0.18cm}
        {\large (ITMO)\par}

        \vspace{0.85cm}

        {\large International Research and Educational Center\par}
        \vspace{0.18cm}
        {\large for Physics of Nanostructures\par}

        \vspace{2.35cm}

        {\Large SYNOPSIS\par}

        \vspace{0.42cm}

        {\large for the subject\par}
        \vspace{0.18cm}
        {\large ``\textit{Research Internship}''\par}

        \vspace{0.56cm}

        {\large on the topic:\par}
        \vspace{0.22cm}
        \makebox[\textwidth][c]{{\large A PERCEPTION-AWARE CONTROL FRAMEWORK}}\par
        \vspace{0.12cm}
        \makebox[\textwidth][c]{{\large FOR AGRICULTURAL QUADROTOR UAVs:}}\par
        \vspace{0.12cm}
        \makebox[\textwidth][c]{{\large DEEP-LEARNING CROP-ANOMALY SEGMENTATION}}\par
        \vspace{0.12cm}
        \makebox[\textwidth][c]{{\large AND TREATMENT MAPPING}}\par
    \end{center}

    \vfill

    \noindent
    \begin{minipage}[t]{0.36\textwidth}
        \raggedright
        {\large\textbf{Student:}\par}
        \vspace{0.18cm}
        {\large\textit{Group No.\ 465672}\par}

        \vspace{0.48cm}

        {\large\textbf{Program:}\par}

        \vspace{0.62cm}

        {\large\textbf{Tutor:}\par}
        \vspace{0.18cm}
        {\large\textit{Associate professor}\par}
    \end{minipage}
    \hfill
    \begin{minipage}[t]{0.46\textwidth}
        \raggedleft
        {\large\textit{Gouda, Ahmed Ashraf}\par}

        \vspace{1.10cm}

        {\large\textit{MSc Mechatronics and Robotics}\par}

        \vspace{0.95cm}

        {\large\textit{Vlasov Sergey Mikhaylovich}\par}
    \end{minipage}

    \vfill

    \begin{center}
        {\large Saint Petersburg -- 2026\par}
    \end{center}
\end{titlepage}
\hypersetup{pageanchor=true}

\tableofcontents
\clearpage

\section*{Abstract}
\addcontentsline{toc}{section}{ABSTRACT}
This work extends a five-controller comparative study for agricultural quadrotor
UAVs (PID, LQR, $H_\infty$, MPC, and SMC) with an opt-in perception-and-decision
layer for precision spraying, and evaluates the five controllers jointly with the
agricultural task they must serve. A downward-facing RGB camera is simulated
during flight; a U-Net (ResNet-18 encoder) semantic segmentation network, trained
on the real \emph{Agriculture-Vision} CVPR-2020 dataset, labels each captured pixel
as background, drydown, weed cluster, or double-plant; affected pixels are
projected into a world-frame treatment map through the camera intrinsics and the
UAV pose; and connected affected zones are revisited by a nearest-neighbour
planner that logs spray events. On the held-out test split the segmentation model
reaches a mean intersection-over-union (mIoU) of $0.405$, which is reported as the
honest, real-imagery accuracy of the system rather than the inflated in-distribution
score of an earlier synthetic prototype. The perception output is scored against a
real, pixel-aligned ground-truth field mask, so that affected-zone recall, coverage,
false-alarm rate, and detection-to-treatment latency measure the controllers on a
genuine detection task. On the agricultural coverage (lawnmower) trajectory, PID,
MPC, and SMC detect every affected zone (recall $1.00$) at $88$--$100\,\%$ coverage,
while LQR and $H_\infty$ leave large gaps and miss half of the affected zones
(recall $0.50$). On the circular stress trajectory every controller detects all
zones, so the discriminating metrics become coverage ($32\%$ for LQR up to $99.5\%$
for MPC) and false-alarm rate. A $10$-seed Monte Carlo campaign with randomized
wind, initial conditions, and sensor noise confirms the ranking is statistically
stable: PID and SMC are the most reliable (tight confidence intervals, perfect
recall), LQR and $H_\infty$ are the weakest and most variable, and MPC attains the
highest coverage but a more variable recall under stress. The study demonstrates
that high survey coverage does not by itself guarantee treatment success, and that
controller choice materially affects the outcome of the full perception-to-action
pipeline.

\noindent\textbf{Keywords:} agricultural UAV, precision agriculture, semantic
segmentation, U-Net, Agriculture-Vision, treatment mapping, perception-aware
control, PID, LQR, $H_\infty$, MPC, sliding mode control.

\clearpage

\section*{Nomenclature}
\addcontentsline{toc}{section}{NOMENCLATURE}

\noindent\textbf{Symbols}
\begin{tabbing}
\hspace{3.2cm}\=\hspace{10cm}\kill
$p=[x,y,z]^{\top}$   \> UAV position in the world (field) frame, m \\
$R\in SO(3)$         \> Body-to-world rotation matrix \\
$v$                  \> UAV linear velocity in the world frame, m/s \\
$\omega$             \> Body angular velocity, rad/s \\
$\phi,\theta,\psi$   \> Roll, pitch, yaw Euler angles, rad \\
$q$                  \> Attitude quaternion, $q=[w,x,y,z]^{\top}$ \\
$m,\,g$              \> Vehicle mass (kg), gravitational acceleration (m/s$^2$) \\
$J$                  \> Body-frame inertia tensor, kg$\cdot$m$^2$ \\
$f,\,\tau$           \> Total thrust (N), body torque vector (N$\cdot$m) \\
$e_p,\,e_v,\,e_R,\,e_\omega$ \> Position, velocity, attitude, angular-rate errors \\
$\Delta t,\,T_{\mathrm{sim}}$ \> Integration step (s), mission horizon (s) \\
$K$                  \> Camera intrinsic matrix \\
$(u,v)$              \> Pixel coordinates in the captured image \\
$H,W$                \> Field length and width, $15\times10$ m \\
$c_{k}$              \> Semantic class of a pixel, $k\in\{0,1,2,3\}$ \\
$\ell(\cdot)$        \> Bayesian log-odds treatment-map accumulator \\
$r$                  \> Treatment-map cell resolution, $0.25$ m \\
$\rho$               \> Spray radius, m \\
$t_{d}$              \> Spray dwell time, s \\
$t_{s}$              \> Spray settling margin, s \\
$\mathrm{mIoU}$      \> Mean intersection-over-union over present classes \\
$R_{\mathrm{aff}}$   \> Affected-zone recall (detected / present) \\
$C$                  \> Coverage of the field footprint, \% \\
$F$                  \> False-alarm rate, m$^{-2}$ \\
$L$                  \> Detection-to-treatment latency, s \\
\end{tabbing}

\noindent\textbf{Acronyms}
\begin{tabbing}
\hspace{3.2cm}\=\hspace{10cm}\kill
UAV     \> Unmanned Aerial Vehicle \\
CNN     \> Convolutional Neural Network \\
IoU     \> Intersection over Union \\
mIoU    \> Mean Intersection over Union \\
ONNX    \> Open Neural Network Exchange \\
GT      \> Ground Truth \\
FPS     \> Frames per Second \\
PID     \> Proportional--Integral--Derivative \\
LQR     \> Linear Quadratic Regulator \\
$H_\infty$ \> H-infinity robust control \\
MPC     \> Model Predictive Control \\
SMC     \> Sliding Mode Control \\
MC      \> Monte Carlo \\
AV      \> Agriculture-Vision (CVPR-2020 dataset) \\
RK4     \> Fourth-order Runge--Kutta integration \\
SO(3)   \> Special Orthogonal group, rotation matrices \\
QP      \> Quadratic Program \\
ARE     \> Algebraic Riccati Equation \\
\end{tabbing}

\clearpage

\section*{Introduction}
\addcontentsline{toc}{section}{INTRODUCTION}
This report extends a comparative study of five trajectory-tracking controllers
for an agricultural quadrotor UAV -- PID, LQR, $H_\infty$, MPC, and SMC -- with a
closed-loop perception and decision layer that turns the flight into a useful
agricultural mission. In the baseline study the controllers were judged only by
how accurately they followed a reference path. Here they are judged by how well the
whole vehicle performs the task of \emph{finding} crop problems from the air and
\emph{planning treatment} for them, because a controller that cannot hold its survey
lines will fail to image parts of the field, and problems that are never imaged can
never be treated. The perception layer consists of a simulated downward camera, a
convolutional segmentation network trained on real aerial farmland imagery, a
projection from image pixels to field coordinates, a probabilistic treatment map,
and a revisit-and-spray planner. The remainder of the report is organized as
follows. Section~\ref{sec:motivation} states the motivation, Section~\ref{sec:related}
reviews related work, and Section~\ref{sec:problem} defines the research problem: the
shared quadrotor platform and its nonlinear dynamics, the constraint and uncertainty
sets, the formulations of the five controllers under evaluation, and the task-level
evaluation metrics. Section~\ref{sec:dl} describes the deep-learning segmentation model
and its accuracy, Section~\ref{sec:pipeline} describes the perception-to-action
pipeline, and Section~\ref{sec:results} reports the simulation results, including the
five-controller comparison and the Monte Carlo campaign. Section~\ref{sec:discussion} discusses the findings,
Section~\ref{sec:limitations} states the limitations and future work, and
Section~\ref{sec:conclusion} concludes.

\clearpage

\section{Motivation}\label{sec:motivation}

\subsection{Scientific Significance}
Autonomous agricultural spraying is a joint perception-and-control problem, yet the
two halves are usually studied separately. Control papers optimize tracking error
against a given reference; perception papers optimize segmentation accuracy on fixed
image datasets. In a real mission the two are coupled: the controller decides
\emph{which} parts of the field the camera ever sees, and the perception model
decides \emph{what} is done about them. A controller that tracks a coverage path
poorly does not merely produce a larger root-mean-square error -- it produces blind
spots in the treatment map, missed infestations, and wasted spray on false
positives. Studying the two together, on the same simulated vehicle and the same
real imagery, exposes trade-offs that neither field can see alone and provides a
more faithful benchmark for controller selection in precision agriculture.

\subsection{Practical Significance (Agricultural UAV Missions)}
Precision spraying aims to apply agro-chemicals only where crop stress, weeds, or
planting anomalies actually occur, reducing cost and environmental load relative to
blanket spraying. This requires the UAV to survey a field, detect affected regions
from the air, and revisit them for targeted treatment. The value of the mission is
therefore measured not by tracking error but by task-level quantities: what fraction
of affected zones are detected (recall), how much of the field is imaged (coverage),
how many false detections trigger unnecessary spray (false-alarm rate), and how long
elapses between detecting a problem and treating it (latency). This report evaluates
the five controllers directly on these quantities.

\subsection{Open Problems and Challenges}
Three challenges motivate the present work. First, segmentation models trained on
synthetic or in-distribution imagery report optimistic accuracy that does not
transfer to real fields; an honest benchmark must use real imagery and a real
ground truth. Second, the perception output must be scored against the \emph{same}
field the vehicle actually flew over, otherwise the detection metrics are
meaningless. Third, high survey coverage is often assumed to imply mission success,
but coverage and treatment precision are different requirements and may rank the
controllers differently. The framework in this report addresses all three.

\section{Related Work}\label{sec:related}
Trajectory-tracking control of quadrotors has been studied extensively. Mahony,
Kumar, and Corke review multirotor modelling and control fundamentals~\cite{mahony2012},
Lee, Leok, and McClamroch derive a geometric tracking controller on $SE(3)$ that is now
a standard nonlinear benchmark~\cite{lee2010}, and Bouabdallah and Siegwart established
an early linear-control comparison on a real platform~\cite{bouabdallah2007}. The five
control families evaluated here rest on classical foundations: optimal linear-quadratic
design~\cite{anderson2007}, $H_\infty$ robust control~\cite{zhou1996}, model predictive
control~\cite{rawlings2017}, and sliding-mode control~\cite{utkin1992}. The companion
controls study to this report compares these five families on agricultural and circular
trajectories under nominal and disturbed conditions; the present work reuses that exact
platform and adds the perception layer on top of it.

Semantic segmentation of imagery advanced rapidly with the U-Net encoder--decoder
architecture~\cite{ronneberger2015} and deep residual encoders~\cite{he2016}, which
together form the ResNet-18 U-Net used here. For aerial farmland specifically, the
\emph{Agriculture-Vision} CVPR-2020 dataset provides large-scale, expert-annotated
imagery of field anomalies such as drydown, weed clusters, and double-plant
regions~\cite{chiu2020}. Precision-agriculture UAV reviews identify accurate
low-altitude tracking under payload and wind variation as the dominant constraint for
spraying missions~\cite{huang2013,zhang2021}, yet they typically treat coverage-path
planning, perception, and low-level control as independent modules.

The contribution here is to close the loop in simulation: to run a real segmentation
model inside the flight simulation of each controller and to score the resulting
treatment decisions against a real ground-truth field, so that the control and
perception layers are evaluated jointly rather than in isolation.

\section{Research Problem Statement}\label{sec:problem}

\subsection{Objectives}
The objectives of this study are: (i) to train and honestly evaluate a real-data
crop-anomaly segmentation model suitable for on-board deployment; (ii) to integrate
that model into the flight simulation of five controllers through a camera model, a
pixel-to-world projection, and a probabilistic treatment map; (iii) to plan and log
targeted spraying over the detected zones; and (iv) to compare the five controllers
on task-level metrics -- affected-zone recall, coverage, false-alarm rate, and
detection-to-treatment latency -- under a single deterministic run and under a
$10$-seed Monte Carlo campaign.

\subsection{System Overview}
The simulated system couples six stages in a closed loop, executed once per camera
frame:
\begin{enumerate}[label=(\arabic*)]
    \item a controller (one of PID, LQR, $H_\infty$, MPC, SMC) flies the quadrotor
          along a reference survey path;
    \item a downward-facing pinhole RGB camera captures the field below the vehicle
          at $5$ frames per second;
    \item the U-Net segmentation model labels every captured pixel with one of the
          four Agriculture-Vision classes;
    \item affected pixels are projected from image coordinates to the world field
          plane using the camera intrinsics $K$ and the UAV pose $(p,R)$;
    \item a Bayesian log-odds treatment map accumulates the projected detections
          into a field-frame grid of $0.25$ m cells;
    \item after coverage, a revisit planner orders the confirmed affected zones and
          a spray executor logs a treatment event for each serviced zone.
\end{enumerate}
Each controller tracks a time-parameterized reference produced by the survey and revisit
planners. The quadrotor state used throughout is
\begin{equation}
    \mathbf{x} = \big[\mathbf{p}^{\top},\,\mathbf{v}^{\top},\,\mathbf{q}^{\top},\,\bm{\omega}^{\top}\big]^{\top} \in \mathbb{R}^{13},
    \label{eq:state}
\end{equation}
comprising position $\mathbf{p}$, linear velocity $\mathbf{v}$, unit attitude quaternion
$\mathbf{q}$, and body angular rate $\bm{\omega}$. The perception layer is evaluated on
top of whichever controller generates the tracking command $\mathbf{x}\!\to\!\mathbf{x}_{\mathrm{ref}}$,
which this study treats as an interchangeable black box, so that the task-level metrics
reflect the controller's effect on the mission rather than its internal control law.

\subsection{Quadrotor Platform and Dynamics}
The perception layer is evaluated on top of the same agricultural quadrotor platform and
the same five controllers as the companion controls study, summarized here so that this
report is self-contained. The vehicle is a rigid body with six degrees of freedom and
four independent inputs, and is therefore an underactuated nonlinear system. Its state is
the position $\mathbf{p}\in\mathbb{R}^3$, linear velocity $\mathbf{v}\in\mathbb{R}^3$,
attitude $R\in SO(3)$, and body angular rate $\bm{\omega}\in\mathbb{R}^3$ of
Eq.~\eqref{eq:state}; the commanded input is the total thrust $f$ and the body torque
$\bm{\tau}$,
\begin{equation}
    \mathbf{u} = [\,f,\ \tau_x,\ \tau_y,\ \tau_z\,]^{\top}. \label{eq:input}
\end{equation}
With $\mathbf{e}_3=[0,0,1]^{\top}$ and gravity $\mathbf{g}$, the continuous-time dynamics
are
\begin{align}
    \dot{\mathbf{p}} &= \mathbf{v}, &
    \dot{\mathbf{v}} &= \mathbf{g} + \tfrac{1}{m}\!\left(fR\mathbf{e}_3 + \mathbf{d}_f\right), \\
    \dot{R} &= R\,\hat{\bm{\omega}}, &
    J\dot{\bm{\omega}} &= \bm{\tau} - \bm{\omega}\times(J\bm{\omega}),
    \label{eq:dyn}
\end{align}
where $\hat{\bm{\omega}}$ is the skew-symmetric matrix of $\bm{\omega}$ and
$\mathbf{d}_f$ is the external disturbance force, primarily wind. Translational motion is
actuated only through attitude, so a controller that tracks its survey lines poorly does
not merely accumulate tracking error -- it leaves parts of the field unimaged and hence
untreatable, which is precisely the coupling this report measures.

The dynamics are advanced at $\Delta t = 0.01$~s ($100$~Hz) by a fourth-order
Runge--Kutta step for $\mathbf{v}$ and $\bm{\omega}$ and a semi-implicit Euler update for
the position, $\mathbf{p}_{k+1} = \mathbf{p}_k + \mathbf{v}_{k+1}\Delta t$. To preserve the geometry of $SO(3)$, the attitude is propagated by the
matrix exponential of the midpoint angular velocity (Rodrigues formula) and re-projected
onto $SO(3)$ by singular-value decomposition,
\begin{equation}
    R_{k+1} = R_k\,\exp\!\big(\hat{\bar{\bm{\omega}}}\,\Delta t\big), \qquad
    R \leftarrow UV^{\top}\ \text{with}\ R = U\Sigma V^{\top}. \label{eq:so3}
\end{equation}
The nominal parameters, held fixed across all five controllers, are mass
$m = 1.3$--$1.5$~kg (empty to loaded), arm length $0.225$~m, inertia
$J = \mathrm{diag}(0.0123,0.0123,0.0224)$~kg$\cdot$m$^2$, per-axis torque limit
$\tau_{\max}=0.5$~N$\cdot$m, over a $15\times10$~m field flown at a spray altitude of
$4$~m and a nominal speed of $\approx 1.2$~m/s. Because the plant is identical for every
controller, any performance difference is attributable to the control law rather than to
the model.

\subsection{Constraints and Uncertainty}
The control inputs are hard-limited to the propulsion capability in every control cycle,
\begin{equation}
    0 \le f \le f_{\max}, \qquad |\tau_i| \le \tau_{\max}, \quad i\in\{x,y,z\}. \label{eq:sat}
\end{equation}
The deterministic runs of Section~\ref{sec:results} use the nominal parameters and a
fixed wind realization. The Monte Carlo campaign repeats both missions over $n=10$
independent seeds; each seed randomizes (i) the wind disturbance $\mathbf{d}_f$, (ii) the
vehicle initial conditions, and (iii) the camera/render and sensor noise seen by the
perception stack. Aggregates are reported as mean $\pm$ a $95\,\%$ confidence interval,
computed as $1.96\,\sigma/\sqrt{n}$ with non-finite runs excluded. This separates a
controller's genuine sensitivity to disturbance and noise from the outcome of a single
favourable seed.

\subsection{The Five Controllers Under Evaluation}\label{sec:controllers}
All five controllers share the state/input interface of
Eqs.~\eqref{eq:state}--\eqref{eq:input} and the actuator limits of Eq.~\eqref{eq:sat};
they differ only in the feedback law $\mathbf{u}=\mathcal{C}_j(\mathbf{x},\mathbf{x}_{\mathrm{ref}})$.
For the linear designs the working state uses the Euler-angle vector
$\bm{\eta}=[\phi,\theta,\psi]^{\top}$. When trajectory feedforward is enabled, a reference
acceleration $\mathbf{a}_{\mathrm{ff}}$ is mapped to a nominal thrust and tilt by
differential flatness, $\mathbf{F}_{\mathrm{ff}} = m(\mathbf{a}_{\mathrm{ff}} + g\mathbf{e}_3)$,
$f_{\mathrm{ff}} = \lVert\mathbf{F}_{\mathrm{ff}}\rVert$,
$\mathbf{b}_{3,d} = \mathbf{F}_{\mathrm{ff}}/\lVert\mathbf{F}_{\mathrm{ff}}\rVert$.

\subsubsection{PID (cascaded)}
An outer translational loop forms the desired force and an inner attitude loop forms the
torque,
\begin{equation}
    \begin{aligned}
        \mathbf{F}_{\mathrm{des}} &= m(g\mathbf{e}_3 + \mathbf{a}_{\mathrm{ff}})
            + K_p^{p}\!\odot\mathbf{e}_p + K_i^{p}\!\odot\!\!\int\!\mathbf{e}_p\,dt
            + K_d^{p}\!\odot\bar{\mathbf{e}}_v, \\
        \bm{\tau} &= K_p^{R}\!\odot\mathbf{e}_R + K_i^{R}\!\odot\!\!\int\!\mathbf{e}_R\,dt
            + K_d^{R}\!\odot\bar{\mathbf{e}}_\omega,
    \end{aligned}
    \label{eq:pid}
\end{equation}
with thrust $f=\mathrm{sat}(\lVert\mathbf{F}_{\mathrm{des}}\rVert, 0.2mg, 1.8mg)$, the
attitude error $\mathbf{e}_R$ formed from the cross product of the current and desired
thrust directions, anti-windup on the integrators, and a $20$~rad/s low-pass filter on
the derivative channels ($\bar{\mathbf{e}}_v,\bar{\mathbf{e}}_\omega$). The implemented
gains are $K_p^{p}=[3.3,3.3,6.5]$, $K_i^{p}=[0.12,0.12,0.25]$, $K_d^{p}=[3.6,3.6,5.8]$,
$K_p^{R}=[7,7,3.5]$, $K_i^{R}=[0.05,0.05,0.02]$, $K_d^{R}=[3,3,1.8]$. PID is simple and
computationally cheap but tuning-sensitive and less systematic about coupling and
disturbances.

\subsubsection{LQR}
A hover/trajectory linearization gives the $12$-state model
$\dot{\mathbf{x}}_\ell = A_c\mathbf{x}_\ell + B_c\mathbf{u}$ with the dominant small-angle
coupling $\dot v_x = g\theta$, $\dot v_y = -g\phi$. The control law adds feedforward
thrust, full-state feedback, and altitude integral action, $\mathbf{u} =
[f_{\mathrm{ff}},0,0,0]^{\top} + K_{\mathrm{LQR}}\mathbf{e} + [mK_{I,z}\!\int e_z\,dt,0,0,0]^{\top}$,
where $K_{\mathrm{LQR}}$ solves the Riccati problem
$\min\int(\mathbf{x}_\ell^{\top}Q\mathbf{x}_\ell + \mathbf{u}^{\top}R\mathbf{u})\,dt$ with
$Q=\mathrm{diag}(2,2,5,\,8,8,12,\,5,5,2,\,3,3,2)$, $R=\mathrm{diag}(2,80,80,100)$, and
altitude integral gain $K_I^z=0.8$. LQR is systematic and smooth but relies on the
linearization and degrades far from hover.

\subsubsection{$H_\infty$}
The same linear model is augmented with a disturbance input,
$\dot{\mathbf{x}}_\ell = A_c\mathbf{x}_\ell + B_c\mathbf{u} + B_w\mathbf{w}$, and the
state-feedback gain minimizes the worst-case disturbance-to-performance gain
$\lVert T_{w\to z}\rVert_\infty < \gamma$ through a bounded-real ARE iteration (with an
LQR fallback), plus a first-order disturbance observer that corrects the velocity-error
block. The design weights are $Q_z=\mathrm{diag}(4,4,10,\,10,10,15,\,6,6,3,\,4,4,2.5)$,
$R_u=\mathrm{diag}(2,60,60,80)$, with $\gamma$ swept over $\{5,6,8,10,15,20\}$; in the
companion controls study the realized closed-loop gain of this same design was certified
a posteriori at $\lVert T_{zw}\rVert_\infty = 2.08 < 5$. $H_\infty$ targets disturbance rejection
explicitly, at the cost of a more involved synthesis.

\subsubsection{Model Predictive Control}
MPC uses the discrete linear model
$\mathbf{x}_{\ell,k+1}=A_d\mathbf{x}_{\ell,k}+B_d\mathbf{u}_k$ over a prediction horizon
$N=20$ and control horizon $N_u=15$, minimizing
\begin{equation}
    \min_{\Delta U}\;
    \begin{aligned}[t]
        &\sum_{i=0}^{N-1}\big(\mathbf{x}_{k+i}-\mathbf{x}^{\mathrm{ref}}_{k+i}\big)^{\top}Q\big(\mathbf{x}_{k+i}-\mathbf{x}^{\mathrm{ref}}_{k+i}\big) \\
        &\quad + \sum_{i=0}^{N-1}\Delta\mathbf{u}_{k+i}^{\top}R\,\Delta\mathbf{u}_{k+i}
            + \big(\mathbf{x}_{k+N}-\mathbf{x}^{\mathrm{ref}}_{k+N}\big)^{\top}Q_N\big(\mathbf{x}_{k+N}-\mathbf{x}^{\mathrm{ref}}_{k+N}\big),
    \end{aligned}
    \label{eq:mpc}
\end{equation}
subject to the model and the input box $u_{\min}\le\Delta u\le u_{\max}$ ($\pm 0.7mg$ on
thrust, $\pm\tau_{\max}$ on torque). The batch QP is solved in closed form
$\Delta U^{\ast}=-H_{\mathrm{qp}}^{-1}f$ and projected onto the box, and only the first
move is applied. The weights are
\begin{equation*}
    Q=\mathrm{diag}(80,80,120,\,30,30,45,\,25,25,8,\,5,5,3),\quad
    R=\mathrm{diag}(0.04,0.18,0.18,0.5),\quad Q_N=20Q.
\end{equation*}
MPC anticipates the reference and respects the constraints directly, but it is the most
computationally demanding law.

\subsubsection{Sliding Mode Control}
SMC defines translational and rotational sliding surfaces
$\mathbf{s}_p=\mathbf{e}_v+\Lambda_p\mathbf{e}_p$ and
$\mathbf{s}_R=\mathbf{e}_\omega+\Lambda_R\mathbf{e}_R$, and commands
\begin{equation}
    \begin{aligned}
        \mathbf{a}_{\mathrm{cmd}} &= \mathbf{a}_{\mathrm{ff}} + \Lambda_p\mathbf{e}_v
            + K_p\,\mathrm{sat}\!\big(\mathbf{s}_p/\phi_p\big), \\
        \bm{\tau} &= J\big(\Lambda_R\mathbf{e}_\omega + K_R\,\mathrm{sat}(\mathbf{s}_R/\phi_R)\big)
            + \bm{\omega}\times(J\bm{\omega}),
    \end{aligned}
    \label{eq:smc}
\end{equation}
with the desired force $\mathbf{F}_{\mathrm{des}}=m(g\mathbf{e}_3+\mathbf{a}_{\mathrm{cmd}})$
and a boundary-layer $\mathrm{sat}(\cdot)$ that replaces the discontinuous sign to reduce
chattering. The surface and gain parameters are $\Lambda_p=\mathrm{diag}(1.6,1.6,2.1)$,
$\Lambda_R=\mathrm{diag}(4,4,2.8)$, $K_p=\mathrm{diag}(1.4,1.4,2.2)$,
$K_R=\mathrm{diag}(3,3,1.8)$, $\phi_p=0.75$, $\phi_R=0.35$. SMC is robust to matched
disturbances but produces more aggressive actuation than the smooth linear laws.

\subsection{Field and Ground Truth}
The vehicle flies over a $15\times10$ m field whose texture and pixel-aligned
ground-truth class mask are both generated from the same real Agriculture-Vision
image patches, so that every affected pixel in the mask corresponds to a real
anomaly in the imagery the camera sees. The field is deliberately built to a
realistic sparsity of $23.1\,\%$ affected area; a denser field would saturate recall
and hide the differences between controllers, while the natural class balance of the
dataset (about $18\,\%$ affected) is preserved in spirit. All detection metrics in
Section~\ref{sec:results} are computed against this real mask.

\subsection{Evaluation Metrics}
The task-level metrics used throughout are defined as follows.
\begin{itemize}
    \item \textbf{Affected-zone recall} $R_{\mathrm{aff}}$: the fraction of
          connected ground-truth affected zones that overlap at least one confidently
          detected treatment cell. This is the primary mission metric, because a
          missed infestation is more costly than a false alarm.
    \item \textbf{Coverage} $C$: the fraction of the field imaged by the camera
          footprint over the survey, in percent.
    \item \textbf{False-alarm rate} $F$: the area of detected affected cells that do
          not correspond to real anomalies, per square metre of field. It must be
          read jointly with coverage, since a vehicle that images more of the field
          also accumulates more false positives.
    \item \textbf{Detection-to-treatment latency} $L$: the mean time between first
          detecting a zone and logging its spray event, in seconds.
    \item \textbf{Segmentation mIoU}: the mean intersection-over-union of the network
          over the present classes, used to report model accuracy independently of
          the flight.
\end{itemize}
Formally, with per-class true positives, false positives, and false negatives
$\mathrm{TP}_k,\mathrm{FP}_k,\mathrm{FN}_k$ (counted in pixels for the network and in map
cells for the flight), the segmentation and task metrics are
\begin{align}
    \mathrm{IoU}_k &= \frac{\mathrm{TP}_k}{\mathrm{TP}_k+\mathrm{FP}_k+\mathrm{FN}_k},
        & \mathrm{mIoU} &= \frac{1}{|\mathcal{C}|}\sum_{k\in\mathcal{C}} \mathrm{IoU}_k,
        \label{eq:iou}\\[2pt]
    R_{\mathrm{aff}} &= \frac{\#\{\text{GT zones hit by a confident cell}\}}{\#\{\text{GT zones}\}},
        & F &= \frac{A_{\mathrm{FP}}}{H\,W}, \label{eq:raff}\\[2pt]
    C &= \frac{A_{\mathrm{imaged}}}{H\,W}\times 100\%,
        & L &= \frac{1}{|\mathcal{Z}|}\sum_{z\in\mathcal{Z}} \big(t^{\mathrm{spray}}_z - t^{\mathrm{det}}_z\big),
        \label{eq:latency}
\end{align}
where $\mathcal{C}$ is the set of classes present, $A_{\mathrm{FP}}$ is the area of
detected affected cells that overlap no real anomaly, $A_{\mathrm{imaged}}$ is the imaged
field area, $H\,W$ is the total field area, and $\mathcal{Z}$ is the set of serviced
zones with detection time $t^{\mathrm{det}}_z$ and spray time $t^{\mathrm{spray}}_z$.

\section{Deep-Learning Segmentation Model}\label{sec:dl}

\subsection{Dataset and Architecture}
The segmentation model is a U-Net with a ResNet-18 encoder, trained on the real
Agriculture-Vision CVPR-2020 dataset (the combined 2017, 2018, and 2019 miniscale
releases). Four classes are used: background (class $0$) and three affected classes
-- drydown ($1$), weed cluster ($2$), and double-plant ($3$). A balanced subset of
$4500$ training, $900$ validation, and $900$ test patches was assembled from the
combined release; the natural class imbalance of the dataset (background dominates
at roughly $82\,\%$ of pixels) makes the affected classes intrinsically difficult.
The network consumes $256\times256$ RGB tiles and is exported to ONNX for CPU
inference inside the MATLAB flight simulation.

\subsection{Training}
The model was trained for $25$ epochs with the Adam optimizer using a resumable
checkpointing scheme. The best validation performance (mean IoU $\approx 0.447$) was
reached at epoch $24$. The training and validation curves are shown in
Figure~\ref{fig:training}.

\figorbox{training_curves.png}
{Training curves of the Agriculture-Vision U-Net.}
{Train loss, validation mean IoU, and per-class validation IoU over $25$ epochs. The
affected classes converge more slowly than background, reflecting their small pixel
share in the dataset.}
{fig:training}

\subsection{Held-out Test Accuracy}
Table~\ref{tab:seg} reports the per-class accuracy on the held-out test split. The
model reaches a mean IoU of $0.405$. This value is far below the $\approx 0.95$ that
an earlier synthetic prototype achieved on its own procedural generator; that
prototype was trained and tested on the same synthetic process and therefore does
not measure field accuracy. It has been removed from the project, and \emph{all}
accuracy claims in this report are the real Agriculture-Vision numbers below.

\begin{table}[H]
    \centering
    \caption{U-Net semantic segmentation accuracy on the held-out Agriculture-Vision
    test split (mean IoU $=0.405$).}
    \label{tab:seg}
    \begin{tabular}{lcccc}
        \toprule
        \textbf{Class} & \textbf{IoU} & \textbf{Precision} & \textbf{Recall} & \textbf{Affected} \\
        \midrule
        background      & 0.812 & 0.896 & 0.896 & no  \\
        drydown         & 0.422 & 0.619 & 0.570 & yes \\
        weed\_cluster   & 0.206 & 0.338 & 0.346 & yes \\
        double\_plant   & 0.180 & 0.300 & 0.310 & yes \\
        \midrule
        \textbf{mean}   & \textbf{0.405} & -- & -- & -- \\
        \bottomrule
    \end{tabular}
\end{table}

Background is segmented well (IoU $0.812$). Among the affected classes drydown is the
strongest (IoU $0.422$, recall $0.570$), weed cluster is moderate (IoU $0.206$), and
double-plant is the weakest (IoU $0.180$). The weakness of double-plant is expected:
it is the rarest and most visually subtle class in the dataset, and it is flagged as
a limitation in Section~\ref{sec:limitations}. The per-class breakdown and the
confusion matrix are shown in Figures~\ref{fig:segbar} and~\ref{fig:segcm}.

\figorbox{segmentation_per_class.png}
{Per-class test accuracy of the segmentation model.}
{Intersection-over-union, precision, and recall for the four classes on the held-out
test split.}
{fig:segbar}

\figorbox{segmentation_confusion.png}
{Row-normalized confusion matrix on the test split.}
{Confusion between the four classes. Most affected-class errors are confusions with
background, consistent with the small pixel share of the anomaly classes.}
{fig:segcm}

\subsection{Qualitative Predictions}
Figure~\ref{fig:qual} shows representative held-out test tiles. Within each tile the
three columns are, from left to right, the input RGB patch, the expert ground-truth
mask, and the model prediction; affected classes are coloured as in the legend (drydown
gold, weed cluster red, double-plant green), and background is left uncoloured. The model
recovers the large drydown and weed-cluster regions well and localizes their boundaries,
but it tends to over-extend the affected area on ambiguous, low-contrast patches and
barely registers the rare double-plant class -- the same behaviour quantified by the
per-class IoU in Table~\ref{tab:seg}.

\begin{figure}[H]
    \centering
    \includegraphics[width=0.34\textwidth]{legend.png}\\[0.5em]
    \begin{subfigure}{0.49\textwidth}
        \centering
        \includegraphics[width=\linewidth]{pred_2017_6Q4R96VAF_2481-3502-2993-4014.png}
        \caption{Drydown region (gold), well recovered.}
    \end{subfigure}\hfill
    \begin{subfigure}{0.49\textwidth}
        \centering
        \includegraphics[width=\linewidth]{pred_2018_ETZD96IZQ_6338-6180-6850-6692.png}
        \caption{Weed cluster (red), close boundary match.}
    \end{subfigure}

    \vspace{0.7em}
    \begin{subfigure}{0.49\textwidth}
        \centering
        \includegraphics[width=\linewidth]{pred_2018_T3QR4RKT1_1161-2370-1673-2882.png}
        \caption{Weed strip along a tree row.}
    \end{subfigure}\hfill
    \begin{subfigure}{0.49\textwidth}
        \centering
        \includegraphics[width=\linewidth]{pred_2018_68XXUAVEB_2181-1137-2693-1649.png}
        \caption{Weed plus a small double-plant (green); prediction over-extends.}
    \end{subfigure}
    \caption{Qualitative segmentation on held-out Agriculture-Vision test tiles.\\\small
    Columns within each panel: input RGB (left), ground truth (middle), model prediction
    (right). The first row shows accurate recovery of drydown and weed cluster; the
    bottom-right panel shows a typical over-prediction on a low-contrast patch and the
    near-invisibility of the rare double-plant class.}
    \label{fig:qual}
\end{figure}
\vspace{0.3em}

\section{Perception-to-Action Pipeline}\label{sec:pipeline}

\subsection{Camera Model and Projection}
A downward-facing pinhole camera with intrinsic matrix $K$ captures the field at the
spray altitude of $4$ m and a frame rate of $5$ FPS. The intrinsics follow from the
sensor field of view,
\begin{equation}
    K =
    \begin{bmatrix}
        f_x & 0   & c_x \\
        0   & f_y & c_y \\
        0   & 0   & 1
    \end{bmatrix},
    \qquad
    f_x = \frac{W/2}{\tan(\theta_h/2)}, \quad
    f_y = \frac{H_{\mathrm{px}}/2}{\tan(\theta_v/2)},
    \label{eq:intrinsics}
\end{equation}
where $\theta_h,\theta_v$ are the horizontal and vertical fields of view,
$(W,H_{\mathrm{px}})$ is the image size in pixels, and $(c_x,c_y)$ is the principal
point. Each detected affected pixel $(u,v)$ is back-projected onto the flat field plane
$z=0$. Writing the pixel in homogeneous form $\tilde{\mathbf{x}}=[u,v,1]^{\top}$, its
viewing ray in the world frame is $\mathbf{r}=R\,K^{-1}\tilde{\mathbf{x}}$, and the
intersection of that ray with the ground plane gives the world coordinate
\begin{equation}
    \begin{bmatrix} x \\ y \end{bmatrix}
    = \big( \mathbf{p} + s\,\mathbf{r} \big)_{1:2},
    \qquad
    s = -\frac{p_z}{r_z},
    \label{eq:backproj}
\end{equation}
where $\mathbf{p}=[x,y,z]^{\top}$ and $R$ are the UAV position and body-to-world
rotation and $r_z$ is the vertical component of $\mathbf{r}$. This maps every affected
pixel into the $15\times10$ m field. The class-identifier convention is kept symmetric
across the Python model (zero-based ids) and the MATLAB pipeline (one-based ids), so
that the affected classes drydown, weed cluster, and double-plant map consistently to
the map-and-spray set.

\subsection{Treatment Map}
Projected detections are accumulated into a field-frame grid of $0.25$ m cells using a
clamped log-odds accumulator. Each time a cell $c$ is observed as affected by a confident
detection, its log-odds are incremented and clamped to a fixed range,
\begin{equation}
    \ell_t(c) = \mathrm{clip}\!\big(\ell_{t-1}(c) + \ell_{\mathrm{hit}},\; \ell_{\min},\, \ell_{\max}\big),
    \label{eq:logodds}
\end{equation}
with $\ell_{\mathrm{hit}}=+0.85$ and $[\ell_{\min},\ell_{\max}]=[-10,10]$; the
implementation accumulates positive evidence only (a negative decrement
$\ell_{\mathrm{miss}}=-0.20$ is reserved in the configuration for future
negative-evidence updates but is not applied in the current pipeline). The posterior
probability that a cell is affected is recovered from its log-odds by the logistic map
\begin{equation}
    P(c\ \text{affected}) = \sigma\big(\ell_t(c)\big) = \frac{1}{1+e^{-\ell_t(c)}},
    \label{eq:sigmoid}
\end{equation}
and a cell enters a treatment zone when $P \ge \tau_p = 0.60$. Since
$\sigma(\ell_{\mathrm{hit}}) = 0.70$, a single confident observation is sufficient to
confirm a cell; the robustness of the map against isolated per-pixel errors therefore
comes from the spatial stage rather than from per-cell temporal integration: connected
affected cells are grouped into zones by connected-component analysis, and only components
of area at least $0.50$ m$^2$ (eight cells) become treatment zones. Scoring the treatment
map against the real ground-truth mask yields the affected-zone recall and false-alarm
rate used in Section~\ref{sec:results}.

\subsection{Revisit and Spray Planning}
After the coverage survey, the confirmed zones are ordered by a nearest-neighbour
heuristic starting from the vehicle's final survey position, and the revisit
reference holds over each zone long enough for the controller to settle within the
spray radius before a spray event is logged. The spray executor is opportunistic:
it services every zone the vehicle passes within the spray radius of, so that a
single hard-to-reach zone does not block the remaining zones. This design replaced
an earlier strictly-sequential executor whose zero settling margin caused
controllers that overshoot on the aggressive circular path to leave zones unsprayed;
the current executor services every detected zone for all five controllers.
Figure~\ref{fig:route} shows an example revisit tour produced by this planner over the
confirmed zones of the real field.

\figorbox{revisit_route.png}
{Planned revisit tour over the detected affected zones.}
{Nearest-neighbour ordering of the confirmed treatment zones (yellow, numbered in
service order) overlaid on the stitched real field. The tour starts from the vehicle's
final survey position; the opportunistic spray executor then services every zone the
vehicle passes within the spray radius of.}
{fig:route}

\section{Simulation Results}\label{sec:results}

\subsection{Real-field Detection}
As a first check that the wired pipeline detects real anomalies, a lawnmower survey
was flown over the real field and the treatment map scored against the ground-truth
mask. The survey imaged $95\,\%$ of the field and recalled every affected zone
($R_{\mathrm{aff}}=1.00$, zone-overlap recall) at a false-alarm rate of
$0.057$ m$^{-2}$, at an on-board inference cost of roughly $80$--$150$ ms per frame
on CPU. The whole-field flyover localizes twelve affected zones -- one drydown patch and
eleven weed clusters -- whose nearest-neighbour revisit tour is shown in
Figure~\ref{fig:route}. Figure~\ref{fig:flyover} compares the real RGB field, the
pixel-aligned ground-truth mask, and the model's whole-field prediction. On the whole stitched field
the fused segmentation reaches mIoU $0.31$ with affected-class precision $0.57$;
these whole-field numbers are lower than the per-patch test mIoU of $0.405$ because
the moving survey windows straddle patch boundaries and see partial, off-nadir views
of each region. This gap is the expected simulation-to-field effect and is reported
openly rather than hidden.

\begin{figure}[H]
    \centering
    \begin{subfigure}{0.32\textwidth}
        \centering
        \includegraphics[width=\linewidth]{field_rgb.png}
        \caption{Real RGB field.}
    \end{subfigure}\hfill
    \begin{subfigure}{0.32\textwidth}
        \centering
        \includegraphics[width=\linewidth]{field_groundtruth.png}
        \caption{Ground-truth mask.}
    \end{subfigure}\hfill
    \begin{subfigure}{0.32\textwidth}
        \centering
        \includegraphics[width=\linewidth]{field_prediction.png}
        \caption{Model prediction.}
    \end{subfigure}
    \caption{Whole-field perception over the real $15\times10$ m field.\\\small The real
    RGB orthomosaic (left), the pixel-aligned ground-truth affected mask (middle), and the
    treatment-map prediction accumulated over the lawnmower survey (right). Affected cells
    are coloured; background is left uncoloured.}
    \label{fig:flyover}
\end{figure}
\vspace{0.3em}

\subsection{Agricultural Trajectory (Coverage Path)}
Table~\ref{tab:lawn} reports the five-controller comparison on the agricultural
lawnmower survey, scored against the real ground truth, for a single deterministic
run.

\begin{table}[H]
    \centering
    \caption{Controller performance on the agricultural lawnmower survey
    (single deterministic run, real ground truth).}
    \label{tab:lawn}
    \begin{tabular}{lcccc}
        \toprule
        \textbf{Controller} & \textbf{Coverage [\%]} & \textbf{Recall} & \textbf{False alarm [m$^{-2}$]} & \textbf{Latency [s]} \\
        \midrule
        PID  & 100.0 & 1.00 & 0.137 & 104.7 \\
        LQR  & 56.4  & 0.50 & 0.147 & 102.5 \\
        Hinf & 62.4  & 0.50 & 0.105 & 91.1  \\
        MPC  & 88.2  & 1.00 & 0.083 & 96.6  \\
        SMC  & 100.0 & 1.00 & 0.123 & 101.9 \\
        \bottomrule
    \end{tabular}
\end{table}

On the coverage mission the controllers separate clearly on the primary metric. PID,
MPC, and SMC track the survey lines tightly enough to image $88$--$100\,\%$ of the
field and detect \emph{every} affected zone (recall $1.00$). LQR and $H_\infty$
follow the lawnmower pattern poorly, image only $56\%$ and $62\%$ of the field, and
consequently miss half of the affected zones (recall $0.50$). Recall tracks coverage:
a controller that cannot follow the survey simply never images some anomalies. The
false-alarm rate is coupled to coverage -- the high-coverage controllers accumulate
slightly more false-positive cells -- and must be read together with coverage rather
than as an independent quality score. Figure~\ref{fig:panel} summarizes the four
metrics across both trajectories.

\figorbox{controller_comparison_panel.png}
{Controller comparison across trajectories (single deterministic run).}
{Coverage, affected-zone recall, false-alarm rate, and mean latency for the five
controllers on the lawnmower and circular trajectories.}
{fig:panel}

\subsection{Circular Trajectory}
Table~\ref{tab:circ} reports the same comparison on the circular stress trajectory.

\begin{table}[H]
    \centering
    \caption{Controller performance on the circular stress trajectory
    (single deterministic run, real ground truth).}
    \label{tab:circ}
    \begin{tabular}{lcccc}
        \toprule
        \textbf{Controller} & \textbf{Coverage [\%]} & \textbf{Recall} & \textbf{False alarm [m$^{-2}$]} & \textbf{Latency [s]} \\
        \midrule
        PID  & 81.0 & 1.00 & 0.061 & 94.9  \\
        LQR  & 32.1 & 1.00 & 0.280 & 145.9 \\
        Hinf & 54.3 & 1.00 & 0.337 & 125.9 \\
        MPC  & 99.5 & 1.00 & 0.158 & 98.4  \\
        SMC  & 76.6 & 1.00 & 0.073 & 105.9 \\
        \bottomrule
    \end{tabular}
\end{table}

On the circular path every controller eventually images enough of the field to detect
all affected zones (recall $1.00$), so recall no longer discriminates. The
discriminating metrics become coverage and false-alarm rate. MPC and PID sweep the
most field ($99.5\%$ and $81.0\%$) at moderate false-alarm rates, whereas LQR and
$H_\infty$ collapse to $32\%$ and $54\%$ coverage and carry by far the highest
false-alarm rates ($0.280$ and $0.337$ m$^{-2}$). With the opportunistic spray
executor every controller now services all detected zones, so the treatment latency
is finite for all five, ranging from $94.9$ s (PID) to $145.9$ s (LQR); the slower,
lower-coverage controllers pay a longer latency because their revisit tours are less
efficient.

\subsection{Monte Carlo Campaign}
To confirm that the ranking is not an artifact of a single seed, a $10$-seed Monte
Carlo campaign was run with randomized wind, initial conditions, and camera noise.
Table~\ref{tab:mc} reports the mean coverage and affected-zone recall with $95\,\%$
confidence intervals.

\begin{table}[H]
    \centering
    \caption{Monte Carlo controller statistics (mean $\pm$ $95\,\%$ CI over $10$ seeds).}
    \label{tab:mc}
    \begin{tabular}{llcc}
        \toprule
        \textbf{Trajectory} & \textbf{Controller} & \textbf{Coverage [\%]} & \textbf{Recall} \\
        \midrule
        agricultural & PID  & $100.0 \pm 0.0$ & $1.00 \pm 0.00$ \\
        agricultural & LQR  & $55.5 \pm 2.5$  & $0.55 \pm 0.10$ \\
        agricultural & Hinf & $62.4 \pm 0.3$  & $0.50 \pm 0.00$ \\
        agricultural & MPC  & $86.6 \pm 0.6$  & $1.00 \pm 0.00$ \\
        agricultural & SMC  & $100.0 \pm 0.0$ & $1.00 \pm 0.00$ \\
        \midrule
        circular & PID  & $81.8 \pm 0.3$   & $1.00 \pm 0.00$ \\
        circular & LQR  & $31.2 \pm 6.6$   & $0.20 \pm 0.22$ \\
        circular & Hinf & $33.9 \pm 11.0$  & $0.75 \pm 0.22$ \\
        circular & MPC  & $90.9 \pm 4.3$   & $0.90 \pm 0.13$ \\
        circular & SMC  & $70.3 \pm 0.8$   & $1.00 \pm 0.00$ \\
        \bottomrule
    \end{tabular}
\end{table}

The Monte Carlo campaign confirms and sharpens the deterministic result. On the
agricultural mission PID and SMC are perfectly reliable (coverage $100\%$, recall
$1.00$, zero confidence interval), MPC is close behind, and LQR and $H_\infty$ are
both weak and, in the LQR case, variable. On the circular mission the separation is
even starker: LQR and $H_\infty$ not only lose coverage but become \emph{unreliable},
with recall dropping to $0.20\pm0.22$ and $0.75\pm0.22$ respectively -- their wide
confidence intervals mean the mission may or may not succeed depending on the
disturbance realization. PID and SMC hold recall at $1.00$ with negligible spread,
and MPC attains the highest coverage but a more variable recall ($0.90\pm0.13$) under
stress. Figure~\ref{fig:mc} shows the campaign graphically.

\figorbox{monte_carlo_panel.png}
{Monte Carlo controller statistics ($n=10$ seeds, $95\,\%$ CI).}
{Coverage and affected-zone recall with confidence bars for the five controllers on
both trajectories. The strong (PID/SMC/MPC) and weak (LQR/$H_\infty$) groups are
separated on coverage.}
{fig:mc}

\subsection{Perception Runtime}
Per-frame ONNX inference runs at roughly $90$--$150$ ms on CPU
(Figure~\ref{fig:runtime}), comfortably within the $5$ FPS camera budget at survey
speed, so perception is not the bottleneck of the pipeline; the limiting factor is
the controller's ability to track the survey path, not the cost of the vision model.

\figorbox{perception_runtime.png}
{Per-frame perception runtime.}
{Mean and $95$th-percentile ONNX inference time per frame for the five controllers.}
{fig:runtime}

\section{Discussion}\label{sec:discussion}
The central finding is that survey coverage and treatment success are related but
not identical requirements, and that controller choice propagates all the way to the
agricultural outcome. On the coverage mission the primary metric, affected-zone
recall, separates the controllers into a strong group (PID, MPC, SMC) that detects
every zone and a weak group (LQR, $H_\infty$) that misses half, purely because the
weak controllers cannot hold the survey lines and therefore never image parts of the
field. On the circular mission recall saturates and the discriminating metrics move
to coverage and false-alarm rate, where the same weak group again trails badly. The
Monte Carlo campaign shows this ranking is statistically robust rather than a
single-seed accident, and it adds a reliability dimension: LQR and $H_\infty$ are not
only worse on average but far more variable under disturbance. A second finding,
exposed during integration, is that high coverage does not automatically produce
successful spraying: an early strictly-sequential spray executor with no settling
margin left zones unsprayed for controllers that overshoot on the aggressive circular
path, which is why the executor was redesigned to be opportunistic and to allow the
controller time to settle. The corrected executor services every detected zone for
all five controllers, and the treatment latency then simply reflects the efficiency
of each controller's revisit tour.

\section{Limitations and Future Work}\label{sec:limitations}
Several limitations should be stated openly. First, the segmentation accuracy is
modest in absolute terms (test mIoU $0.405$), which is honest for real aerial
farmland but far below the inflated scores obtainable on synthetic data; the
double-plant class in particular is weak (IoU $0.180$ on the test split and near zero
on the moving survey) and would benefit from more training data or a larger
Agriculture-Vision-2021 model on a GPU. Second, the whole-field survey metrics are
lower than the per-patch test metrics because the moving camera sees partial,
off-nadir views; a higher frame rate or overlap-aware fusion would narrow this gap.
Third, the camera and field models are simplified -- a flat field plane, a real
orthomosaic texture, and idealized lighting -- so field-deployment claims require
on-board validation flights. Fourth, the treatment latency depends on the revisit-tour
ordering rather than on the control law alone, so it should be read as a mission-planning
cost rather than a pure controller metric; the Monte Carlo campaign was re-run with the
corrected opportunistic spray executor, so its latency statistics are now consistent with
the deterministic latencies in Tables~\ref{tab:lawn} and~\ref{tab:circ}, with every seed
servicing all detected zones (which tightens the latency confidence intervals relative to
the earlier strictly-sequential executor). None of these limitations affect the coverage
and recall rankings, which are the principal result.

\section{Conclusion}\label{sec:conclusion}
This report integrated a real-data crop-anomaly segmentation model into the flight
simulation of five controllers and evaluated them jointly on a genuine
perception-to-action task. Trained on the Agriculture-Vision CVPR-2020 dataset, the
U-Net reaches an honest held-out test mIoU of $0.405$, and its detections are scored
against a real, pixel-aligned field mask. On the agricultural coverage mission PID,
MPC, and SMC detect every affected zone at high coverage, while LQR and $H_\infty$
miss half; on the circular stress mission the same weak group collapses in coverage
and false-alarm rate; and a $10$-seed Monte Carlo campaign confirms that PID and SMC
are the most reliable controllers, LQR and $H_\infty$ the weakest and most variable,
and MPC the highest-coverage but more variable option under stress. The study shows
that a perception-aware evaluation ranks the controllers on the quantity that
actually matters for precision agriculture -- how many crop problems are found and
treated -- and that this ranking is driven by each controller's ability to hold its
survey path.

\begin{thebibliography}{99}
\addcontentsline{toc}{section}{REFERENCES}

\bibitem{mahony2012}
R. Mahony, V. Kumar, and P. Corke, ``Multirotor aerial vehicles: modeling, estimation,
and control of quadrotor,'' \emph{IEEE Robotics \& Automation Magazine}, vol.~19,
no.~3, pp.~20--32, 2012.

\bibitem{lee2010}
T. Lee, M. Leok, and N. H. McClamroch, ``Geometric tracking control of a quadrotor UAV
on $SE(3)$,'' in \emph{Proc. IEEE Conf. on Decision and Control (CDC)}, 2010,
pp.~5420--5425.

\bibitem{bouabdallah2007}
S. Bouabdallah and R. Siegwart, ``Full control of a quadrotor,'' in \emph{Proc.
IEEE/RSJ Int. Conf. on Intelligent Robots and Systems (IROS)}, 2007, pp.~153--158.

\bibitem{anderson2007}
B. D. O. Anderson and J. B. Moore, \emph{Optimal Control: Linear Quadratic Methods},
Dover, 2007.

\bibitem{zhou1996}
K. Zhou, J. Doyle, and K. Glover, \emph{Robust and Optimal Control}, Prentice Hall,
1996.

\bibitem{rawlings2017}
J. B. Rawlings, D. Q. Mayne, and M. Diehl, \emph{Model Predictive Control: Theory,
Computation, and Design}, Nob Hill, 2017.

\bibitem{utkin1992}
V. I. Utkin, \emph{Sliding Modes in Control and Optimization}, Springer, 1992.

\bibitem{ronneberger2015}
O. Ronneberger, P. Fischer, and T. Brox, ``U-Net: Convolutional networks for biomedical
image segmentation,'' in \emph{Proc. Int. Conf. on Medical Image Computing and
Computer-Assisted Intervention (MICCAI)}, 2015, pp.~234--241.

\bibitem{he2016}
K. He, X. Zhang, S. Ren, and J. Sun, ``Deep residual learning for image recognition,''
in \emph{Proc. IEEE Conf. on Computer Vision and Pattern Recognition (CVPR)}, 2016,
pp.~770--778.

\bibitem{chiu2020}
M. T. Chiu et al., ``Agriculture-Vision: A large aerial image database for agricultural
pattern analysis,'' in \emph{Proc. IEEE/CVF Conf. on Computer Vision and Pattern
Recognition (CVPR)}, 2020, pp.~2828--2838.

\bibitem{huang2013}
Y. Huang, W. C. Hoffmann, Y. Lan, W. Wu, and B. K. Fritz, ``Development of a spray
system for an unmanned aerial vehicle platform,'' \emph{Applied Engineering in
Agriculture}, vol.~25, no.~6, pp.~803--809, 2013.

\bibitem{zhang2021}
Y. Zhang, D. Yuan, Y. Wei, and J. Liu, ``A review of unmanned aerial vehicle application
in precision agriculture: concepts, tools, and techniques,'' \emph{Remote Sensing},
vol.~13, no.~6, 1204, 2021.

\end{thebibliography}

\end{document}
